\subsection*{Drug representations do not improve within-drug ranking}

We evaluated all major drug representation types under ridge regression with 10-fold drug-blind CV on PASO splits. No encoding improved per-drug $r$ beyond $\Delta = +0.001$ of the cell-only baseline (Fig.~\ref{fig:decomp}b, Extended Data Table~1). This null held for Morgan fingerprints, ChemBERTa embeddings, drug-target vectors, and LINCS L1000 perturbation signatures\cite{subramanian2017next}. The LINCS result is particularly informative because these signatures encode a drug's transcriptional effect rather than its chemical structure, ruling out the hypothesis that structural representations fail because they lack mechanistic content. The null also held under scaffold-stratified splits ($\Delta = +0.001$).

An 11.7M-parameter Transformer encoder replicated the null under the same protocol. Morgan fingerprints gave $\Delta = +0.008$, and LINCS and drug-target vectors gave no significant improvement. An omics ablation confirmed that RNA-seq alone captures nearly all predictive signal, with mutations and further modalities adding negligibly (Fig.~\ref{fig:decomp}c). Cell-representation expressiveness is not the bottleneck.

The null replicated across three independent datasets (CTRPv2, BeatAML, and PRISM, Extended Data Table~3). On PRISM, the effect was marginal, consistent with its pharmacologically diverse viability-based assay. The uniformity of the null across all encoding types shifts the question from representation quality to training distribution. Does MoA information fail because it cannot be encoded, or because the training distribution is misaligned?
